%% CV - Njome David (Version Francaise - MASTER / mine d'or, pas pour export one-pager direct)
%% Compiler : xelatex ou pdflatex
\documentclass[10pt, a4paper]{article}

% -- Packages -----------------------------------------------------------------
\usepackage[a4paper, top=1.2cm, bottom=1.2cm, left=1.4cm, right=1.4cm]{geometry}
\usepackage{fontenc}
\usepackage[utf8]{inputenc}
\usepackage{microtype}
\usepackage{parskip}
\usepackage{titlesec}
\usepackage{enumitem}
\usepackage{tabularx}
\usepackage{array}
\usepackage{xcolor}
\usepackage{hyperref}
\usepackage{multirow}
\usepackage{booktabs}

% -- Colors ---------------------------------------------------------------------
\definecolor{navyblue}{HTML}{1F4E79}
\definecolor{darkgray}{HTML}{595959}
\definecolor{lightgray}{HTML}{F2F2F2}

% -- Hyperlinks -------------------------------------------------------------------
\hypersetup{
  colorlinks=true,
  urlcolor=navyblue,
  linkcolor=navyblue
}

% -- No page number -----------------------------------------------------------------
\pagestyle{empty}

% -- Section style ------------------------------------------------------------------
\titleformat{\section}
  {\large\bfseries\color{navyblue}}
  {}{0em}
  {}
  [\vspace{-6pt}\textcolor{navyblue}{\rule{\linewidth}{1.2pt}}\vspace{2pt}]

\titlespacing{\section}{0pt}{10pt}{6pt}

% -- List style -----------------------------------------------------------------------
\setlist[itemize]{
  leftmargin=1.2em,
  itemsep=1pt,
  parsep=0pt,
  topsep=2pt,
  label={\textcolor{navyblue}{\textbullet}}
}

% -- Entry command: Titre \hfill Meta ------------------------------------------------
\newcommand{\entryhead}[2]{%
  \noindent\textbf{#1}\hfill{\small\textit{\textcolor{darkgray}{#2}}}\\[-4pt]%
}

\newcommand{\entrysub}[1]{%
  {\small\textit{\textcolor{darkgray}{#1}}}\\[2pt]%
}

% -- Label en gras -----------------------------------------------------------------------
\newcommand{\blabel}[1]{\textbf{#1}\ }

% ===============================================================================
\begin{document}
% ===============================================================================

% -- EN-TETE ---------------------------------------------------------------------------
{\LARGE\bfseries\color{navyblue} NJOME DAVID}\\[4pt]
{\small\textcolor{darkgray}{%
  Melen, Yaoundé\quad|\quad
  \href{mailto:njomedavid26@gmail.com}{njomedavid26@gmail.com}\quad|\quad
  +237 671 711 462\quad|\quad
  \href{https://www.linkedin.com/in/njome-david-002a11213}{linkedin.com/in/njome-david-002a11213}\quad|\quad
  \href{https://github.com/Njome-David}{github.com/Njome-David}
}}

\vspace{4pt}
\textcolor{navyblue}{\rule{\linewidth}{1.8pt}}
\vspace{2pt}

% -- PROFIL PERSONNEL -----------------------------------------------------------------
\section{Profil Personnel}

Étudiant en quatrième année de Génie Informatique à l'École Nationale Supérieure Polytechnique
de Yaoundé (ENSPY), avec une base pratique en développement full-stack, machine learning et IA, programmation système,
réseaux et systèmes embarqués/IoT, renforcée récemment par un stage à
orientation recherche sur la communication IoT hors-réseau. Entré à l'ENSPY directement depuis
un enseignement secondaire anglophone vers un cursus majoritairement francophone, avec une
progression académique constante depuis. Aspire à évoluer vers l'Intelligence Artificielle et le
Cloud Computing, avec un intérêt à long terme pour l'application de l'IA aux enjeux de santé et
d'agriculture en contexte africain, dans une démarche attentive à l'éthique de l'IA. Candidat
pour approfondir cette base technique dans un cadre académique exigeant.

% -- FORMATION -------------------------------------------------------------------------
\section{Formation}

\entryhead{École Nationale Supérieure Polytechnique de Yaoundé (ENSPY) - Génie Informatique}{09/2023 -- Présent}
\entrysub{Yaoundé, Cameroun}
\begin{itemize}
  \item \blabel{Niveau :} 4\textsuperscript{e} année. 3\textsuperscript{e} année validée avec
        \textbf{Mention Très Bien}, 1\textsuperscript{er} de la promotion ; MGP cumulée
        \textbf{3,30} sur 3 ans (MGP niveau 3~: 3,34).
  \item \blabel{IA et Fondements Formels :} Systèmes Formels et IA (logique des propositions et
        prédicats, Prolog) ; Probabilités et Statistiques ; Analyse Numérique (interpolation
        polynomiale, intégration/dérivation numérique, régression) ; Bases de Données (SQL).
  \item \blabel{Systèmes et Réseaux :} Systèmes d'Exploitation (gestion des processus, de la
        mémoire et des fichiers) ; Programmation Système (C, appels système) ; Architecture des
        Ordinateurs (portes logiques, pseudo-assembleur, tableaux de Karnaugh) ; Introduction aux
        Réseaux (approche théorique complète, nombreux travaux pratiques Cisco, niveau CCNA
        complet).
  \item \blabel{Algorithmique :} Structures de données et algorithmes fondamentaux, analyse de
        complexité (Grand~O), programmation dynamique.
  \item \blabel{Conception des Systèmes d'Information :} Modélisation Merise (MCD/MLD) et UML,
        normalisation de bases de données.
  \item \blabel{Également :} Mathématiques de Base (fondements de la théorie de la mesure),
        Anglais, Management.
\end{itemize}

\vspace{4pt}
\entryhead{Collège Bilingue de Laval - Enseignement Secondaire Général}{09/2016 -- 07/2023}
\entrysub{Douala, Cameroun}
\begin{itemize}
  \item Classé \textbf{8\textsuperscript{e} à l'échelle nationale} au GCE A-Level -
        Mathematics~A, Physics~A, ICT~A, Further Maths~A, Chemistry~A
\end{itemize}

% -- EXPERIENCE ----------------------------------------------------------------------
\section{Expérience}

\entryhead{Laboratoire IoT de l'ENSPY \normalfont\small- Stagiaire}{07/2026 -- 09/2026}
\entrysub{Yaoundé, Cameroun - stage en équipe, sous la responsabilité de Dr Anne Marie Chana (Chef de laboratoire)}
\begin{itemize}
  \item Contribution à un projet d'équipe visant un dispositif de surveillance de pièce
        (température, humidité, concentration de gaz, détection de présence) avec tableau de
        bord temps réel, conçu pour évoluer vers un système multi-dispositifs couvrant un
        logement entier.
  \item Phase MVP : capteurs sur ESP32 avec liaison WiFi/HTTP vers une base de données ;
        objectif final : fonctionnement entièrement hors-réseau via \textbf{LoRaWAN} sur une
        carte Wyres sous \textbf{RIOT OS}.
  \item \blabel{Responsabilité personnelle :} l'ensemble de la chaîne de communication
        hors-réseau - réception des données via les passerelles LoRaWAN, exploitation d'un
        serveur local, décodage, parsing et stockage des données en base pour alimenter le
        tableau de bord.
  \item \blabel{Savoir-être démontré :} adaptation à un cadre de recherche ouvert et
        expérimental, dont l'objectif technique a évolué en cours de projet (du MVP WiFi vers
        une solution entièrement hors-réseau en LoRaWAN) ; autonomie de travail au sein d'une
        petite équipe de recherche tout en restant aligné avec les objectifs plus larges du
        laboratoire ; communication technique régulière avec la chef de laboratoire ; résolution
        de problèmes d'intégration matériel-logiciel sans procédure toute faite ; rigueur de
        documentation, avec la production d'un rapport de stage et d'un rapport de projet.
  \item Exposition concrète à LoRa/LoRaWAN et RIOT OS pour la communication hors-réseau -
        directement pertinent pour étendre la connectivité en zones rurales, un axe de
        recherche actif du laboratoire. Première expérience de travail en structure de
        recherche.
\end{itemize}

% -- PROJETS ----------------------------------------------------------------------
\section{Projets Universitaires / Personnels}

\entryhead{Application CROSS Authentication \normalfont\small- Projet de Groupe (Chef de projet)}{2025}
\begin{itemize}
  \item Plateforme d'authentification de documents officiels par cryptographie post-quantique
        (schéma de signature CROSS).
  \item \blabel{Compétences mobilisées :} Gestion de projet, documentation technique, Python.
  \item Pilotage d'une équipe de 6 personnes sur 7 semaines, de la conception à la livraison
        d'une plateforme post-quantique pleinement fonctionnelle.
\end{itemize}

\vspace{4pt}
\entryhead{Salle d'Impression Illustrée \normalfont\small- Projet Individuel}{2025}
\begin{itemize}
  \item Simulation système d'une salle d'impression concurrente où des utilisateurs aux
        priorités, volumes et types d'impression variés se disputent des imprimantes limitées -
        illustrant les principes de concurrence et de synchronisation.
  \item \blabel{Compétences mobilisées :} Concurrence et synchronisation, conception
        d'algorithmes d'ordonnancement, analyse de performance, programmation C.
  \item Conception et comparaison expérimentale d'un ordonnanceur par priorité original face
        aux algorithmes FIFO et SJF - avec de meilleurs résultats d'équité, mesurés sur le
        temps d'attente, le temps de rotation et le taux de famine.
\end{itemize}

\vspace{4pt}
\entryhead{Application de Gestion de Production (App Suite) \normalfont\small- Collaboration}{2025 -- 2026}
\begin{itemize}
  \item Réalisation du frontend complet d'une application multi-tenant de gestion de production
        pour entreprises.
  \item \blabel{Compétences mobilisées :} Architecture frontend, conception d'interfaces
        multi-tenant, développement full-stack collaboratif.
  \item Intégrée à une suite applicative plus large maintenue et déployée par le Pr Djotio
        Thomas pour des professionnels africains ; propriété du dépôt transférée à ce dernier
        pour la maintenance continue.
\end{itemize}

\vspace{4pt}
\entryhead{DIGISCHOOL \normalfont\small- Projet d'Équipe (En cours)}{2025 -- Présent}
\begin{itemize}
  \item Application web de gestion d'école primaire - utilisateurs, élèves et parents.
  \item \blabel{Compétences mobilisées :} React, Node.js, modélisation UML/cas d'utilisation,
        développement agile en équipe.
  \item Construite à partir de son propre travail de modélisation UML sur la gestion de la
        discipline scolaire bilingue ; projet en cours.
\end{itemize}

\vspace{4pt}
\entryhead{FairShare \normalfont\small- Projet de Hackathon}{2025 -- 2026}
\begin{itemize}
  \item Application de partage de dépenses réalisée en conditions de hackathon.
  \item \blabel{Compétences mobilisées :} Next.js 14, Prisma, Supabase, Framer Motion,
        prototypage rapide sous contrainte de temps.
  \item Version stable, déployée et téléchargeable en tant que PWA (Progressive Web App).
\end{itemize}

\vspace{4pt}
\entryhead{Levons-nous et Bâtissons - Plateforme de Sondage \normalfont\small- Projet Individuel}{2025}
\begin{itemize}
  \item Application web de sondage sur mesure pour une initiative civique chrétienne destinée
        aux étudiants de l'ENSPY.
  \item \blabel{Compétences mobilisées :} HTML/CSS/JS, Google Apps Script, design UX itératif.
  \item Réalisée à travers trois itérations de design vers une identité visuelle camerounaise
        sur mesure, avec un backend écrivant vers Google Sheets.
\end{itemize}

\vspace{4pt}
\entryhead{Gestion des Cotisations d'Organisation \normalfont\small- Collaboration Club GI}{2025}
\begin{itemize}
  \item Backend Laravel pour la gestion des adhésions et cotisations d'une organisation.
  \item \blabel{Compétences mobilisées :} Laravel, authentification Sanctum, tests d'API
        (Postman).
  \item Modules Auth et CRUD Membre livrés avec une suite de 13 tests Postman couvrant le cas
        nominal, la sécurité et la validation.
\end{itemize}

\vspace{4pt}
\entryhead{GeoAccel \normalfont\small- Outil Personnel}{2025}
\begin{itemize}
  \item Application de bureau Windows calculant l'accélération maximale du sol via quatre lois
        d'atténuation historiques (McGuire, Petrovski, Joyner-Boore, Sabetta-Pugliese).
  \item \blabel{Compétences mobilisées :} Python, Tkinter, méthodes numériques, modélisation de
        domaine technique.
  \item Réalisée pour un camarade de Génie Civil afin de remplacer un calcul manuel fastidieux,
        livrée comme outil autonome fonctionnel.
\end{itemize}

\vspace{4pt}
\entryhead{RT-Comops - Service STOCKMAN \normalfont\small- Contribution}{2025}
\begin{itemize}
  \item Conception d'un service de gestion de stock pour un noyau monolithique modulaire en
        architecture hexagonale.
  \item \blabel{Compétences mobilisées :} PlantUML, modélisation de domaine, architecture
        hexagonale, rédaction technique.
  \item Introduction d'un discriminant counterpartyType et d'un agrégat StockMovementLink pour
        corriger un défaut de modélisation dans la conception PartyRef existante.
\end{itemize}

\vspace{4pt}
\entryhead{Tarification Étudiante pour Plateformes IA - Étude de Marché \normalfont\small- Projet Individuel}{2025}
\begin{itemize}
  \item Étude de marché entrepreneuriale plaidant pour une tarification étudiante officielle
        des plateformes IA au Cameroun.
  \item \blabel{Compétences mobilisées :} Conception de sondages, étude de marché, ReportLab
        (génération de PDF), Python.
  \item Étude complète produite à partir de 100 réponses à un sondage, ciblant Anthropic comme
        plateforme prioritaire et le paiement mobile-money comme exigence non négociable.
\end{itemize}

% -- COMPETENCES ----------------------------------------------------------------
\section{Compétences}

\vspace{2pt}
\noindent
\begin{tabularx}{\linewidth}{@{} >{\bfseries}p{4.8cm} X @{}}
  Programmation
    & Python, Java, C/C++, HTML/CSS, JavaScript, GitHub, \LaTeX \\[3pt]
  Développement Full-Stack
    & Next.js, React, Node.js, Laravel, Spring Boot, PostgreSQL, MySQL \\[3pt]
  Systèmes \& Réseaux
    & Fonctionnement OS, appels système en C, architecture des ordinateurs, réseaux Cisco/CCNA \\[3pt]
  Données, IA \& Maths Appliquées
    & SQL, Python (Pandas, bases du Machine Learning), logique formelle et Prolog, méthodes
      numériques (interpolation, régression), probabilités et statistiques \\[3pt]
  Conception des Systèmes
    & Merise, UML, normalisation de bases de données, analyse des besoins \\[3pt]
  Leadership \& Gestion de Projet
    & Coordination d'équipe, livraison par sprints, chef de groupe récurrent \\[3pt]
  Communication
    & Documentation technique, prise de parole en public, reporting transverse \\
\end{tabularx}

% -- LANGUES ----------------------------------------------------------------
\section{Langues}

\vspace{2pt}
\noindent
\begin{tabularx}{\linewidth}{@{} >{\bfseries}p{4.8cm} X @{}}
  Français  & Courant \\[3pt]
  Anglais   & Courant \\[3pt]
  Espagnol  & Notions \\[3pt]
  Allemand  & Notions \\
\end{tabularx}

% -- CERTIFICATIONS ET DISTINCTIONS ----------------------------------------------
\section{Certifications et Distinctions}

\begin{itemize}
  \item Certification en \textbf{Microsoft Excel 2016}, classé \textbf{1\textsuperscript{er}
        au Cameroun} et qualifié pour représenter le pays au \textbf{Microsoft Office Specialist
        (MOS) World Championship} (2021 et 2022).
  \item \textbf{Séminaire MAPLE} - NASEY OPTICA Chapter (Optica / OSA), ENSPY :\\
        \textit{2\textsuperscript{e} Éd.} (fév.--avr. 2025) : \textbf{Prix du Meilleur Étudiant}.\\
        \textit{3\textsuperscript{e} Éd.} (fév.--avr. 2026) : \textbf{Formateur Assistant et
        Organisateur}.
\end{itemize}

% -- ACTIVITES EXTRASCOLAIRES ----------------------------------------------------
\section{Activités Extrascolaires}

\begin{itemize}
  \item \textbf{Député Junior du Cameroun auprès des Nations Unies} - représentation de la
        jeunesse nationale à travers des engagements diplomatiques, prises de parole en public
        et débats structurés sur les ODD.
  \item Membre actif de la cellule projets du \textbf{Club GI ENSPY}.
  \item Passionné de basketball et joueur de jeux vidéo en équipe - sens de la stratégie,
        de la collaboration et de la compétition.
\end{itemize}

% -- REFERENCES ----------------------------------------------------------------
\section{Références}

\textit{\textcolor{darkgray}{Disponibles sur demande.}}

\end{document}
